% ============================================================================
%  Chapitre 5 — Mise en œuvre
%  Cible : 20 à 25 pages. C'est le cœur du mémoire. À rédiger J7-J9.
%
%  MÉTHODE pour chaque section : (1) objectif et exigence couverte,
%  (2) conception retenue, (3) extrait de code commenté, (4) difficulté
%  rencontrée et solution. Le point (4) est ce qui distingue un mémoire
%  d'ingénieur d'une documentation technique : ne pas le sauter.
% ============================================================================
\chapter{Mise en œuvre}
\label{chap:mise-en-oeuvre}

\section{Vue d'ensemble}
\label{sec:impl-vue-ensemble}

\aecrire{Une page. Rappeler les trois chaînes et présenter le tableau
récapitulatif ci-dessous, puis annoncer le plan du chapitre.}

\begin{table}[H]
\centering
\caption{Chaînes d'intégration mises en œuvre}
\label{tab:impl-workflows}
\begin{tabularx}{\textwidth}{@{}l X >{\raggedright\arraybackslash}p{3.6cm}@{}}
\toprule
\textbf{Chaîne} & \textbf{Objet} & \textbf{Déclenchement} \\
\midrule
\texttt{pm77xx-app} & Construction, analyse et test de la couche applicative &
  Planifié (jours ouvrés), manuel, ou appelé \\
\texttt{pm77xx-yocto} & Construction des variantes d'image et du \gls{sdk} &
  Planifié (jours ouvrés) ou manuel \\
\texttt{publish-image} & Construction, validation et publication de l'image de
  compilation & Manuel \\
\bottomrule
\end{tabularx}
\end{table}

\chiffre{Durée moyenne d'exécution de chacune des trois chaînes, à ajouter en
colonne supplémentaire une fois les mesures collectées (chapitre~\ref{chap:evaluation}).}

% ===========================================================================
\section{Environnement de compilation conteneurisé}
\label{sec:impl-conteneur}

\emph{Exigences couvertes : EF1, ENF1. Sous-question : QR1.}

\subsection{Objectif}
\aecrire{Rappeler la limite L1 : divergence des environnements entre postes
de développement et entre postes et intégration. Objectif : une image de
conteneur unique, versionnée, servant à la fois au développeur et à la chaîne
d'intégration.}

\subsection{Composition de l'image}

L'image est construite à partir d'une base Ubuntu~24.04 \gls{lts} durcie et
maintenue en interne. L'ensemble des outils y est installé à version épinglée,
de manière à ce que deux constructions séparées dans le temps utilisent une
chaîne d'outils strictement identique.

\begin{table}[H]
\centering
\caption{Principaux composants de l'image de compilation et versions épinglées}
\label{tab:impl-image-versions}
\begin{tabularx}{\textwidth}{@{}l X l@{}}
\toprule
\textbf{Composant} & \textbf{Rôle} & \textbf{Version} \\
\midrule
GCC / G++ & Compilation C et C++ & 13.3.0 \\
CMake & Système de construction applicatif & 3.28.3 \\
Meson / Ninja & Construction de la brique de sécurité & 1.3.2 / 1.11.1 \\
Coverity Analysis & Analyse statique & 2024.12.0 \\
GoogleTest & Cadre de test unitaire & 1.12.1 \\
CMocka & Simulation de dépendances en test & 1.1.7 \\
gcovr & Mesure de couverture de code & 7.0 \\
Doxygen / Graphviz & Génération de documentation & 1.9.8 / 2.42.2 \\
\bottomrule
\end{tabularx}
\end{table}

\aecrire{Compléter avec les bibliothèques applicatives (cJSON, libsystemd,
libwebsockets, libevent, libmodbus, SQLite, RapidJSON) : un tableau ou une
énumération courte, pas les deux.}

\subsection{Points de conception notables}

\aecrire{Développer les points suivants, qui montrent une réflexion
d'ingénierie et non une simple recopie de commandes.}

\begin{description}[style=nextline,leftmargin=0.8cm]
  \item[Construction multi-étapes] Les composants compilés depuis les sources
    sont produits dans une étape intermédiaire puis copiés dans l'image
    finale, ce qui évite d'y conserver les sources et les outils de
    construction.
  \item[Accès aux dépôts privés pendant la construction] La récupération de la
    brique de cybersécurité nécessite une authentification. Un montage de
    secret temporaire est utilisé, de sorte que la clé n'apparaisse dans
    aucune couche de l'image finale. \aecrire{Développer : c'est un point de
    sécurité concret, le jury peut interroger dessus.}
  \item[Utilisateur non privilégié] L'image définit un utilisateur dédié
    plutôt que d'exécuter en tant que superutilisateur.
    \aecrire{Signaler honnêtement l'écart : l'exécution en intégration se fait
    malgré tout en superutilisateur pour cause de compatibilité avec le
    service d'exécuteurs. Expliquer pourquoi, et le reporter en perspectives.}
\end{description}

\subsection{Cycle de vie de l'image}
\aecrire{Décrire la chaîne \texttt{publish-image} : construction, analyse du
Dockerfile par un analyseur dédié, vérification de la présence et des versions
des outils, validation par exécution de la chaîne applicative complète contre
l'image candidate, puis publication sous une étiquette définitive. Insister
sur le fait qu'une image n'est publiée qu'après avoir prouvé qu'elle
construit le produit : c'est le point de conception intéressant.}

\begin{figure}[H]
\centering
\aecrire{Figure : cycle de vie de l'image de compilation (construction,
validation, publication).}
\caption{Cycle de validation et de publication de l'image de compilation}
\label{fig:impl-cycle-image}
\end{figure}

% ===========================================================================
\section{Chaîne de construction applicative}
\label{sec:impl-pipeline-app}

\emph{Exigences couvertes : EF2, EF6, EF7, ENF5. Sous-question : QR1.}

\subsection{Déclencheurs et paramètres}
\aecrire{Trois modes de déclenchement : planifié en jours ouvrés, manuel avec
paramètres, et appelable depuis une autre chaîne. Expliquer pourquoi le mode
appelable est indispensable à la validation de l'image de compilation
(section précédente).}

\begin{table}[H]
\centering
\caption{Paramètres d'entrée de la chaîne applicative}
\label{tab:impl-parametres-app}
\begin{tabularx}{\textwidth}{@{}l X l@{}}
\toprule
\textbf{Paramètre} & \textbf{Rôle} & \textbf{Défaut} \\
\midrule
\texttt{app\_repo\_branch} & Branche du dépôt applicatif & \texttt{develop} \\
\texttt{test\_repo\_branch} & Branche du dépôt de tests & \texttt{main} \\
\texttt{builder\_img\_tag} & Étiquette de l'image de compilation & épinglée \\
\texttt{ci\_repo\_branch} & Branche du dépôt d'intégration & \texttt{main} \\
\texttt{run\_cov\_analysis} & Activation de l'analyse statique & \texttt{true} \\
\bottomrule
\end{tabularx}
\end{table}

\subsection{Matrice de construction}
\label{ssec:impl-matrice}

La construction est déclinée selon deux axes : le type de construction
(\texttt{DEBUG}, \texttt{RELEASE}) et l'architecture cible (\texttt{amd64} pour
la compilation native de validation, \texttt{aarch64} pour la compilation
croisée destinée au produit), soit quatre combinaisons exécutées en parallèle.

\aecrire{Justifier pourquoi certaines étapes ne sont exécutées que sur une
branche de la matrice : les tests unitaires et d'intégration s'exécutent sur
\texttt{DEBUG}/\texttt{amd64} (instrumentation de couverture et exécution
native possible), l'analyse de composition porte sur
\texttt{RELEASE}/\texttt{aarch64} (artefact effectivement livré). Cette
répartition est un arbitrage coût/couverture : l'expliciter.}

\begin{table}[H]
\centering
\caption{Répartition des activités sur la matrice de construction}
\label{tab:impl-matrice-activites}
\begin{tabular}{@{}lcccc@{}}
\toprule
\textbf{Activité} & \textbf{DEBUG} & \textbf{DEBUG} & \textbf{RELEASE} & \textbf{RELEASE} \\
 & \textbf{amd64} & \textbf{aarch64} & \textbf{amd64} & \textbf{aarch64} \\
\midrule
Compilation            & \checkmark & \checkmark & \checkmark & \checkmark \\
Analyse statique       & \checkmark & \checkmark & --- & --- \\
Tests unitaires        & \checkmark & --- & --- & --- \\
Couverture de code     & \checkmark & --- & --- & --- \\
Tests d'intégration    & \checkmark & --- & --- & --- \\
Analyse de composition & --- & --- & --- & \checkmark \\
\bottomrule
\end{tabular}
\end{table}

\subsection{Séquence d'exécution}
\begin{figure}[H]
\centering
\includegraphics[width=0.85\textwidth]{pm77xx-app-workflow}
\caption{Diagramme d'activité de la chaîne de construction applicative}
\label{fig:impl-activite-app}
\end{figure}

\aecrire{Commenter la figure en suivant le flux d'exécution. Une figure non
commentée dans le texte est considérée comme absente.}

\subsection{Gestion des dépôts multiples et des sous-modules}
\aecrire{Difficulté concrète à développer : le code applicatif est réparti sur
un dépôt principal et de nombreux sous-modules privés. Expliquer le mécanisme
d'authentification retenu par clé de service et son incidence sur la durée de
récupération des sources.}

\begin{lstlisting}[language=yaml,caption={Extrait : récupération des sources
  et exécution dans le conteneur de compilation},label={lst:impl-checkout}]
# TODO : coller ici l'extrait réel du workflow, réduit aux lignes utiles.
# Règle : un listing ne doit jamais dépasser une demi-page. Au-delà,
# le mettre en annexe A et n'en citer que la référence.
\end{lstlisting}

\subsection{Difficultés rencontrées}
\aecrire{Les documenter explicitement. Pistes déjà identifiées dans la
documentation existante : nettoyage des services en fin d'exécution supposant
des répertoires absents sur un exécuteur neuf, gestion du délai d'attente au
démarrage des services, exclusion des tests destinés au cœur temps réel qui
ne peuvent pas s'exécuter sur l'hôte.}

% ===========================================================================
\section{Chaîne de construction de la distribution embarquée}
\label{sec:impl-pipeline-yocto}

\emph{Exigence couverte : EF3. Sous-question : QR1.}

\subsection{Organisation de la configuration}
La configuration de construction est décrite de manière déclarative avec kas et
structurée en trois niveaux : la carte cible, le projet d'image, et les options
de construction.

\aecrire{Décrire les niveaux : fichier de carte définissant la machine
\texttt{\Carte}, fichiers de projet définissant la distribution et la cible
d'image, fichiers d'option pour la génération du \gls{sdk} et les réglages
d'intégration.}

\subsection{Variantes d'image}
\begin{table}[H]
\centering
\caption{Variantes d'image produites}
\label{tab:impl-variantes}
\begin{tabularx}{\textwidth}{@{}l X@{}}
\toprule
\textbf{Variante} & \textbf{Usage} \\
\midrule
\texttt{dev} & Image de développement, outils de débogage inclus \\
\texttt{hwtest} & Validation matérielle en production \\
\texttt{netboot} & Démarrage par le réseau (initramfs) \\
\texttt{rescue} & Image de secours (initramfs) \\
\texttt{prod} & Image de production, noyau optimisé pour le temps de démarrage \\
\bottomrule
\end{tabularx}
\end{table}

\subsection{Maîtrise du temps de construction}
\label{ssec:impl-cache}

\aecrire{Section à fort contenu technique, à soigner. Expliquer le mécanisme
du cache d'état partagé et du cache de téléchargement, leur mutualisation
entre exécutions, et le réglage du parallélisme. Fournir la mesure comparative
avec et sans cache : c'est l'un des résultats les plus parlants du
chapitre~6.}

\chiffre{Durée de construction à froid et à chaud, par variante d'image.}

\subsection{Génération et publication du kit de développement}
\aecrire{Expliquer que le \gls{sdk} généré pour la variante de développement
est publié dans le référentiel d'artefacts, ce qui permet aux développeurs
applicatifs de compiler pour la cible sans reconstruire la distribution.
C'est un apport concret à l'équipe : le mettre en valeur.}

\subsection{Construction des paquets du banc de test}
\aecrire{Décrire la construction séparée des paquets destinés à
l'orchestrateur du banc de test matériel, sur une base distincte de la cible
embarquée.}

% ===========================================================================
\section{Qualité logicielle : tests et couverture}
\label{sec:impl-qualite}

\emph{Exigences couvertes : EF6, EF7, EF9.}

\subsection{Tests unitaires}
\aecrire{Décrire l'exécution des tests unitaires via le lanceur de CMake,
l'exclusion des tests destinés au cœur temps réel, et l'instrumentation de
couverture activée en construction de débogage.}

\subsection{Tests d'intégration}
Les tests d'intégration sont écrits avec Robot Framework et organisés en suites
correspondant aux domaines fonctionnels : gestion administrative et système,
gestion des équipements, gestion des défauts, notifications, et interfaces nord
Modbus, BACnet et \gls{edm}.

\aecrire{Décrire la séquence : démarrage des services applicatifs dans le
conteneur, attente de disponibilité, exécution des suites, arrêt des services,
collecte des rapports. Souligner la difficulté principale : exécuter un
ensemble de services système dans un conteneur sans gestionnaire de services
complet.}

\subsection{Restitution des résultats}
\aecrire{Décrire la génération automatique d'une synthèse en Markdown injectée
dans le rapport d'exécution : indicateur visuel de taux de réussite, tableau
des compteurs, liste des tests en échec. Même principe pour la couverture, avec
des seuils configurables. Argument à défendre : un résultat qui n'est pas lu
n'a aucune valeur ; la restitution fait partie de la conception, pas de la
cosmétique.}

\begin{figure}[H]
\centering
\aecrire{Capture d'écran de la synthèse d'exécution. Sources disponibles :
\texttt{ci/docs/images/UTReports.png}, \texttt{ci/docs/images/buildArtifacts.png}.}
\caption{Synthèse automatique des résultats de test publiée à chaque exécution}
\label{fig:impl-synthese-tests}
\end{figure}

% ===========================================================================
\section{Sécurité : analyse statique et analyse de composition}
\label{sec:impl-securite}

\emph{Exigences couvertes : EF4, EF5. Sous-question : QR2.}

\subsection{Analyse statique}
L'analyse statique est réalisée en trois phases distinctes : la capture, qui
instrumente la construction et enregistre les unités de compilation ; l'analyse
proprement dite, qui applique les vérificateurs ; et la remontée des résultats
vers le serveur d'analyse.

\aecrire{Justifier le découpage en trois phases et son intérêt en intégration
continue. Détailler la configuration : activation de l'ensemble des
vérificateurs, y compris les vérificateurs de sécurité de la famille C,
résolution des appels virtuels et des pointeurs de fonction, normalisation des
chemins pour que les défauts restent corrélables d'une exécution à l'autre
malgré des répertoires de travail éphémères. Ce dernier point est une
difficulté réelle propre aux exécuteurs jetables : la mettre en avant.}

\begin{lstlisting}[language=yaml,caption={Extrait de la configuration
  d'analyse statique},label={lst:impl-cov-config}]
# TODO : coller l'extrait réel (phases capture / analyze / commit).
\end{lstlisting}

\subsubsection{Restitution ciblée des défauts}
\aecrire{Décrire le filtrage par vues de priorité 1 et 2, restreintes au code
écrit par l'équipe à l'exclusion des composants tiers. Expliquer le client
d'extraction développé : interrogation de l'interface de programmation du
serveur, attente de fin d'analyse, mise en forme des résultats en rapport
consultable. Justifier le développement de cet outil plutôt que l'usage de
l'interface web : l'objectif est que le développeur voie les défauts sans
quitter son rapport d'exécution.}

\subsubsection{Politique de blocage}
\aecrire{Point de discussion à traiter honnêtement. Situation actuelle : la
détection de défauts n'interrompt pas l'exécution, seule une défaillance de
construction ou de connexion au serveur le fait. Justifier ce choix par la
dette initiale, puis proposer la trajectoire cible : blocage sur les défauts
nouvellement introduits uniquement. Cette lucidité sera mieux notée qu'une
présentation idéalisée.}

\subsection{Analyse de composition des artefacts}
\aecrire{Décrire la soumission de l'artefact de production à l'analyseur de
composants binaires, l'attente du résultat, la récupération du rapport et
l'extraction des indicateurs : nombre de composants identifiés, nombre de
composants porteurs de vulnérabilités, répartition des vulnérabilités par
sévérité. Relier explicitement aux obligations de nomenclature logicielle
présentées en section~\ref{sec:art-supply-chain}.}

\begin{table}[H]
\centering
\caption{Indicateurs extraits de l'analyse de composition}
\label{tab:impl-bdba-metriques}
\begin{tabularx}{\textwidth}{@{}l X@{}}
\toprule
\textbf{Indicateur} & \textbf{Définition} \\
\midrule
Composants identifiés & Nombre total de composants tiers détectés dans l'artefact \\
Composants vulnérables & Composants pour lesquels au moins une vulnérabilité est connue \\
Vulnérabilités par sévérité & Répartition critique / élevée / moyenne / faible \\
\bottomrule
\end{tabularx}
\end{table}

\subsection{Gestion des secrets}
\aecrire{Recenser les secrets nécessaires (clé de service Git, jetons
d'authentification aux serveurs d'analyse, jeton du référentiel d'artefacts) et
décrire leur mode d'injection ainsi que les précautions prises pour éviter
toute divulgation dans les journaux d'exécution. Couvre ENF4.}

% ===========================================================================
\section{Infrastructure d'exécution}
\label{sec:impl-runners}

\emph{Exigence couverte : ENF3. Sous-question : QR4.}

\subsection{Modèle d'exécuteurs éphémères}
\aecrire{Décrire le service mutualisé d'exécuteurs : allocation d'une machine
neuve par exécution, destruction en fin d'exécution. Détailler les
conséquences de conception : aucun cache local persistant, durée de
récupération des sources et de l'image à chaque exécution, chemins de travail
variables. Relier au problème de normalisation des chemins évoqué en
section~\ref{sec:impl-securite}.}

\subsection{Description de l'infrastructure en code}
\aecrire{Décrire les éléments décrits en Terraform : réservation d'adresses
réseau fixes pour les besoins de filtrage vers les serveurs internes, rôles
d'accès au fournisseur cloud, politiques d'autorisation par dépôt. Décrire le
processus de modification : proposition de changement, validation automatique
des politiques, application après fusion.}

\begin{lstlisting}[language=hcl,caption={Extrait de la description
  d'infrastructure},label={lst:impl-terraform}]
# TODO : coller un extrait représentatif, anonymisé.
\end{lstlisting}

\subsection{Considérations de sécurité}
\aecrire{Traiter : isolation entre exécutions, portée minimale des jetons,
absence de persistance, risque lié à l'exécution en superutilisateur signalé
en section~\ref{sec:impl-conteneur}.}

% ===========================================================================
\section{Observabilité de la chaîne}
\label{sec:impl-observabilite}

\emph{Exigences couvertes : EF9, ENF6. Sous-question : QR3.}

\subsection{Motivation}
\aecrire{Une synthèse par exécution ne permet pas de percevoir une tendance.
L'objectif est de conserver l'historique des indicateurs pour piloter
l'évolution de la qualité, et non seulement constater l'état d'une exécution.}

\subsection{Chaîne de collecte}
\begin{figure}[H]
\centering
\aecrire{Figure : chaîne exécution $\rightarrow$ exportateur $\rightarrow$
base de séries temporelles $\rightarrow$ tableaux de bord.}
\caption{Chaîne de collecte et de restitution des indicateurs}
\label{fig:impl-observabilite}
\end{figure}

\aecrire{Décrire la production des indicateurs au format d'import de la base
de séries temporelles, leur envoi, et leur restitution sous forme de tableaux
de bord. Préciser le déploiement : conteneurs orchestrés, chart de
déploiement, gestion chiffrée des secrets, exposition avec certificat
d'entreprise.}

\subsection{Indicateurs retenus}
\begin{table}[H]
\centering
\caption{Indicateurs collectés et historisés}
\label{tab:impl-indicateurs}
\begin{tabularx}{\textwidth}{@{}l X l@{}}
\toprule
\textbf{Famille} & \textbf{Indicateur} & \textbf{Source} \\
\midrule
Construction & Durée, issue, fréquence & Plateforme d'orchestration \\
Test & Tests exécutés, réussis, en échec & Rapports de test \\
Couverture & Lignes, fonctions, branches & Outil de couverture \\
Analyse statique & Défauts de priorité 1 et 2 & Serveur d'analyse \\
Composition & Composants et vulnérabilités par sévérité & Analyseur binaire \\
\bottomrule
\end{tabularx}
\end{table}

\section{Synthèse du chapitre}
\label{sec:impl-synthese}

\aecrire{Reprendre le tableau des exigences du chapitre~4 et indiquer pour
chacune la section qui la couvre et son état : couverte, partiellement
couverte, non couverte. Assumer les exigences non couvertes, elles alimentent
les perspectives.}
